\section{Conclusions, Limitations, and Future Work}
\label{sec:conclusions}

This letter defined the SAVNav task, where an embodied agent navigates to an acoustic target while treating social sound events as spatial social boundaries. To address this, we proposed SAVMap, a spatio-temporal representation that converts transient audio observations into persistent spatial constraints, and the SAVNav policy, which integrates active sensory exploration with socially-compliant navigation. Evaluations in both simulation and on a physical robot demonstrate that our approach reduces NLOS collisions and improves PSC over vision-based baselines, underscoring the value of persistent acoustic-to-spatial memory for safe navigation.

The current system presents two main limitations. First, it relies on a pre-computed static semantic and topological map to predict human emergence and anchor acoustic observations. Second, the acoustic simulation assumes a relatively simple environment without strong interfering noise, which may not fully reflect the acoustic complexity of real-world scenarios.

Future work will extend the framework to unknown environments by incorporating simultaneous topological and acoustic mapping. Additionally, building on our preliminary real-robot deployment, we plan to scale up physical experiments and evaluate the system under unconstrained real-world social dynamics.
